\section{The whole degree of the actual regular pulled-line trace}

\begin{lemma}[Native canonical-line unit and the same regular pulled-line trace]
\label{lem:kn-canonical-regular-pulled-line-native-unit-degree}
\lean{EnvelopeProperty.Vanishing.NadelGlobalCanonicalTrace.unit_comp_regularPulledCanonicalTrace}
\leanok
Keep the original smooth base and tame SNC auxiliary data, the entire
original cover $\pi:Y\to X$, its original whole degree $D$, and the
regular pulled-line trace $R:\pi_*\pi^*\omega_X\to\omega_X$ constructed
from the actual global structure trace by the original projection formula.
Then the native adjunction unit satisfies $\eta_{\omega_X}R=D\,1_{\omega_X}$.
This trace has source $\pi_*\pi^*\omega_X$; a regular trace on $\omega_Y$
is a separate assertion.
\uses{def:nadel-regular-pulled-canonical-trace,lem:line-pf-native-adjunction-unit,thm:kn-structure-trace-global-unit-degree}
\end{lemma}
\begin{proof}
Use the proved actual projection-formula normalization of the native
adjunction unit. Cancel the original pushed unit constraint with its
inverse, followed by the original projection formula with its inverse.
The remaining map is the actual global structural-unit/structure-trace
composite tensored with the original canonical line. Apply the already
proved global whole-degree law. Invertible tensoring is an additive
equivalence, so it preserves the natural-number multiple of the identity.
Conjugation by the same original unit constraint gives the result.
No generic base change, field trace interface or incoming degree equation
is used for this global identity.
\end{proof}
